%\ifodd\value{page}\hbox{}\newpage\fi
\chapter*{Śrī Lalitā Sahasranāma — Śloka 5}
\addcontentsline{toc}{chapter}{Śloka 5 - Nomi 10,11}

\begin{sloka}
{\sanskritfont
\begin{tabular}{c}
अष्टमीचन्द्रविभ्राजदलिकस्थलशोभिता ।\\
मुखचन्द्रकलङ्काभमृगनाभिविशेषका ॥ ५ ॥
\end{tabular}

\vspace{1.5em}

{\middlesize \iastfont
aṣṭamīcandra-vibhrāja-dalika-sthala-śobhitā |\\
mukha-candra-kalaṅkābha-mṛganābhi-viśeṣakā ||5|| \\}

\vspace{1em}

{\middlesize
\anvaya{%
अष्टमीचन्द्रविभ्राजदलिकस्थलशोभिता
मुखचन्द्रकलङ्काभमृगनाभिविशेषका ।
}
}

\middlesize \iastfont \itmean{%
«Adornata sulla fronte dallo splendore della luna dell’ottavo giorno;  
il cui volto, simile alla luna, è impreziosito da un segno di muschio che richiama una macchia lunare.»%
}


}
\end{sloka}

\wordmeaning{%
\wordline{अष्टमी-चन्द्र-विभ्राज-दलिक-स्थल-शोभिता}{aṣṭamī-candra-vibhrāja-dalika-\\
sthala-śobhitā}
  { adornata sulla fronte dallo splendore della luna dell’ottavo giorno;}%
\wordline{मुख-चन्द्र-कलङ्क-आभ-मृगनाभि-विशेषका}{mukha-candra-kalaṅka-ābha-\\
mṛganābhi-viśeṣakā}
  { il cui volto lunare è impreziosito da un segno di muschio che richiama una macchia sulla luna;}%
}

Lo \textit{Śloka} 5 prosegue la descrizione della forma di \textit{Lalitā} soffermandosi sul volto e sulla regione della fronte. L’immagine della luna dell’ottavo giorno (\textit{aṣṭamīcandra}) non suggerisce una pienezza compiuta, ma una luce giovane e misurata, che esprime grazia contenuta ed equilibrio. La menzione del \textit{dalika-sthala} indica che è la curva stessa della fronte a essere paragonata a questo crescente, in consonanza con quelle tradizioni esegetiche che descrivono la fronte di \textit{Lalitā} come un semicerchio rovesciato quando emerge dal fuoco sacrificale. Il volto, paragonato alla luna, è inoltre segnato da un tilaka di muschio che richiama la macchia lunare, a indicare che ciò che potrebbe apparire come imperfezione diventa ornamento nel corpo divino: la bellezza non consiste nell’assenza di segni, ma nella loro integrazione armoniosa, nel manifestarsi – nel punto del “terzo occhio” – di una luminosità calma, neutra e imperturbata che riflette lo stato interiore di \textit{Lalitā}.

Letteralmente, \textit{mṛga-nābhi} è un composto di \textit{mṛga} («cervo», in particolare il cervo muschiato) e \textit{nābhi} («ombelico»), e designa l’“ombelico del cervo”, sede del sacco ghiandolare da cui si ricava il muschio naturale (\textit{kastūrī}), celebre per il suo profumo penetrante. Sul piano simbolico, \textit{mṛga-nābhi} evoca il \textit{tilaka} di muschio tracciato sulla fronte all’altezza dell’\textit{ājñā-cakra} come segno di concentrazione, protezione e risveglio interiore. Nel \textit{nāma} \textit{mukha-candra-kalaṅkābha-mṛganābhi-viśeṣakā}, il volto-luna di \textit{Lalitā} è segnato da questo punto di muschio, ornamento iniziatico, traccia profumata di una coscienza raccolta e interiorizzata.

Lo \textit{Śloka} 5 presenta due nāma effettivi. \textbf{\textit{Aṣṭamīcandra-vibhrāja-dalika-sthala-śobhitā}} designa \textit{Lalitā} come colei la cui fronte è adornata dalla luce della luna all’ottavo \textit{tithi}, simbolo di misura e controllo della potenza luminosa, in cui la luna si presenta come arco equilibrato piuttosto che come pienezza esplosiva. \textbf{\textit{Mukha-candra-kalaṅkābha-mṛganābhi-viśeṣakā}} afferma che il volto stesso, simile alla luna, è impreziosito da un segno di muschio, mostrando come la perfezione divina includa e trasfiguri il segno senza negarlo, trasformando la “macchia” in tilaka che abbellisce il viso e lo carica di significato iniziatico.
​